\chapter{引言}
\label{ch:Introduction}
广义标量形式的非线性薛定谔方程（NLSE）描述了复电场归一化包络 $\A=\A(z,T)$ 在存在损耗、色散以及 $\chi^3$ 非线性的介质中传播时的演化规律。对应的复电场为 $\E=\E(z,T)=\E_0\cdot \exp(-i(\betag(\omega_0)z-\omega_0T))$，其中 $\omega_0$ 是载波角频率，$\betag(\omega_0)$ 是载波空间频率。其数学形式为

\begin{align}
    \label{eq:GNLSE}
    \partial_z \A = \frac{\alpha}{2}\A+i \sum_{n=2}^{\infty}i^n \frac{\betag_n}{n!}\partial_T^n\A  + i\gamma\left(1+\frac{i}{\omega_0}\partial_T  \right)\left(
\A \int_{0}^{\infty} R(T_{delay})|\A(z,T-T_{delay})|^2 dT_{delay} \right),
\end{align}
其中 $\alpha$ 是功率损耗/增益系数，$\betag_n=\partial_\omega^n\betag(\omega)|_{\omega=\omega_0}$ 是在 $\omega=\omega_0$ 处对空间频率做泰勒展开得到的系数，$\gamma$ 是非线性参数，$R(T_{delay})$ 是描述当前时刻 $T$ 之前延迟 $T_{delay}$ 处非线性响应的时间响应函数。求解式~\ref{eq:GNLSE} 可以用来描述超连续谱产生~\cite{supercontinuum_original_paper,NLSE_original}、孤子~\cite{soliton_first_theory,Soliton_experimental_first}、光纤通信系统中的非线性噪声~\cite{poggiolini2014detailedanalyticalderivationgn}，以及许多具有丰富科学与工业应用的特殊光学现象。

\section{目标}
本教程按式~\ref{eq:GNLSE} 的各个组成项及其相互作用展开，目标是帮助读者建立对相关数学与底层物理的直观理解。为此，文中有意略去一些更复杂的效应，例如与光偏振相关的内容，转而更详细地推导和举例说明纯标量情形下的各种现象。希望这种处理方式能够为读者提供分析非线性光学常见实验结果的基本工具，并为进一步阅读更高级的论文、教材和其他资料打下基础。

\section{获取方式}
本教程可在 \href{https://github.com/OleKrarup123/NLSE-primer}{GitHub} 上免费获取，作者也会根据读者反馈持续更新。欢迎将问题、建议和勘误发送至 \href{yourfavouriteta@gmail.com}{YourFavouriteTA@gmail.com}。


\section{引用说明}
本教程既引用了所讨论主题的原始学术文献，也链接到 YouTube 视频、个人网页、在线百科、交互式工具等材料。这些内容既包括业余爱好者制作的直观讲解，也包括专业研究者提供的高质量资源。这样做的目的是一方面为读者提供系统梳理正式文献的起点，便于今后独立撰写非线性光学方向的论文或学位论文；另一方面也收集了大量易于理解、质量较高的辅助解释，帮助读者加深理解。

\subsection{关于引用本教程}
本教程不包含关于 NLSE 的原创研究，更适合作为一份详细的学习笔记。在撰写介绍新研究成果的论文或学位论文时，不应把它当作原始资料来引用。更合适的做法是直接引用相关主题最早的原创工作，例如讨论孤子时，应优先引用~\cite{soliton_first_theory} 和~\cite{Soliton_experimental_first}。不过，如果你的目的只是帮助读者先熟悉 NLSE，从而更容易理解原始研究，那么可以按如下 BibTeX 格式引用本教程~\cite{NLSE_primer}：\\
\begin{boxA}
@misc\{NLSE\_primer,  \\
author = "O. Krarup", \\
title = "\{A Primer on the Nonlinear Schrödinger Equation\}", \\
note = "Commit SHA: 9f1b93a", \\
url = \{https://github.com/OleKrarup123/NLSE-primer/blob/main/NLSE\_primer.pdf\}\}
\end{boxA}

需要注意的是，由于本教程托管在 GitHub 上且会持续更新，URL 应始终指向最新版本，而 “Commit SHA” 则应填写实际引用时对应的最新提交哈希值。
